% Point to the main file
\documentclass[../../Main.tex]{subfiles}

% ========== Start the document ==========
\begin{document}

\section{Section 1.7 Intro to Proofs}

What do we prove? 
Usually that some pattern or formula is true.

\begin{example}
    \begin{align*}
        6 + 8 &= 14 \\
        12 + 16 &= 28 \\
        2 + 2 &= 4 \\
    \end{align*}

    We believe that the sum of 2 even integers is even.

    Try to prove that: even + even = even

    \medskip

    \begin{work}
        \begin{center}
            An even number can be defined as $2n$, where n is an integer
            $$2n + 2m = 2(n + m)$$
        \end{center}
    \end{work}

    \begin{proof}
        \begin{center}
            let $a, b$ be even integers.

            Then there exists integers $n, m$ such that $a = 2n, b=2m$

            $$ a + b = 2n + 2m = 2(m + n) $$ is even, as desired
        \end{center}
    \end{proof}

\end{example}

But writing it this way is too long. Let's define some notation:

\subsection{Set Notation}

\begin{tablewithheader}
{|c|c|c|}
{ Name & Symbol & Set Notation }
    
    Set of all natural numbers & $\mathbb{N}$ & $\{ 0, 1, 2, 3, \ldots \}$ \\
    Set of all integers & $\mathbb{Z}$ & $\{ 0, -1, 1, -2, 2, \ldots \}$ \\
    Set of positive integers & $\mathbb{Z}^+$ & $\{ 1, 2, 3, \ldots \}$ \\
    Set of negative integers & $\mathbb{Z}^-$ & $\{ -1, -2, -3, \ldots \}$ \\
    Set of rational numbers & $\mathbb{Q}$ & $\{ \frac{a}{b} \mid a,b \in \mathbb{Z}, b\not= 0 \}$ \\
    Set of real numbers & $\mathbb{R}$ & Limit of all convergent sequences of real numbers \\
    Set of all complex numbers & $\mathbb{C}$ & $\{ a+bi \mid a,b \in \mathbb{R}, i^2 =-1 \}$ \\

\end{tablewithheader}

$\exists$ means there exists

Let's rewrite the previous example using this new notation

\begin{center}
    $$a,b \in \mathbb{Z} \text{, both even}$$
    Then, $\exists m, n \in \mathbb{Z} \text{ such that}$
    $$a = 2m, b = 2n$$

    Their sum $a + b = 2m + 2n = 2(m + n)$

    is even, as desired
\end{center}

% continue writing the other examples. For now, I'm going to skip them since I'm lazy

\subsection{Methods of Proving $p \implies q$}

\begin{enumerate}
    \item
        \textbf{Direct Proof:}

        Start with $p$ and $math$ then conclude $q$

    \item
        \textbf{Indirect Proof:}

        Prove the contrapositive proof $\neg q \implies \neg p$

    \item
        \textbf{Finite Cases:}

        Use the fact that $p \equiv p_1 \lor p_2 \lor \ldots \lor p_n$

        Then prove $(p_1 \implies q) \land (p_2 \implies q) \land \ldots \land (p_n \implies q)$

\end{enumerate}

% ========== End the document ==========
\end{document}